\documentclass[10pt,reqno,a4paper]{amsart}

\usepackage[utf8]{inputenc}
\usepackage[english]{babel}
\usepackage{amsmath,amssymb,amsthm,amsfonts}
\usepackage{mathtools}
\usepackage[margin=25mm]{geometry}
\usepackage{booktabs}
\usepackage{array}
\usepackage{seqsplit}
\usepackage{microtype}
\usepackage[colorlinks=true,linkcolor=blue,urlcolor=blue]{hyperref}
\usepackage[capitalize]{cleveref}

\numberwithin{equation}{section}

\theoremstyle{plain}
\newtheorem{theorem}{Theorem}[section]
\newtheorem{proposition}[theorem]{Proposition}
\newtheorem{lemma}[theorem]{Lemma}
\newtheorem{corollary}[theorem]{Corollary}
\theoremstyle{definition}
\newtheorem{definition}[theorem]{Definition}
\theoremstyle{remark}
\newtheorem{remark}[theorem]{Remark}

\newcommand{\R}{\mathbb{R}}
\newcommand{\N}{\mathbb{N}}
\newcommand{\Pp}{\mathbb{P}}
\newcommand{\Ee}{\mathbb{E}}
\newcommand{\1}{\mathbf{1}}
\newcommand{\abs}[1]{\left|#1\right|}
\newcommand{\norm}[1]{\left\lVert#1\right\rVert}
\newcommand{\ind}{\Rightarrow}
\DeclareMathOperator{\Var}{Var}

% Declaration names are long and TeX cannot break inside \texttt. Turn every
% underscore into a break opportunity, with no hyphen inserted.
\newcommand{\ubreak}{\char`\_\discretionary{}{}{}}
\newcommand{\lean}[1]{{\ttfamily\let\_\ubreak\seqsplit{#1}}}

\title[First-passage central limit theorem]{The positive-drift first-passage
central limit theorem}
\author{Benny Avelin}
\date{\today}

\begin{document}

\begin{abstract}
  We prove the first-passage central limit theorem used in the fixed-width
  cutoff argument.  The proof uses only the ordinary central limit theorem,
  the strong law of large numbers, finite-dimensional invariance of an
  i.i.d. sequence under reversal, Slutsky's theorem, and Portmanteau.  The
  main estimate is that the terminal drawdown from the running maximum is
  uniformly tight.  We also prove the moving-level and time-perturbation
  forms required by the paper and transfer the result to a random walk with
  an increasing almost surely bounded correction.  The result is used in
  \S3 of \emph{Absorption cutoff and stationary singularities for rounded
  Gaussian random dynamical systems}, where it is cited from the literature;
  the proof given here is the one carried out in the accompanying Lean
  development.
\end{abstract}

\maketitle
\enlargethispage{2pt}

\section{Statement and notation}\label{sec:statement}

Let $(\Omega,\mathcal F,\Pp)$ be a probability space, and let
$(\xi_j)_{j\geq0}$ be independent identically distributed real random
variables such that
\begin{equation}\label{eq:moments}
  \xi_0\in L^2(\Omega,\mathcal F,\Pp),
  \qquad
  \mu:=\Ee\xi_0>0,
  \qquad
  \sigma:=\sqrt{\Var(\xi_0)}>0.
\end{equation}
Thus $\sigma^2=\Var(\xi_0)$.  We state square-integrability separately
because positivity of the variance is not used as a substitute for the
moment hypothesis.
Put
\begin{equation}\label{eq:walk-maximum-drawdown}
  S_0:=0,
  \qquad
  S_n:=\sum_{j=0}^{n-1}\xi_j,
  \qquad
  M_n:=\max_{0\leq k\leq n}S_k,
  \qquad
  D_n:=M_n-S_n.
\end{equation}
For $u>0$, define the first weak ascending passage time
\begin{equation}\label{eq:hitting-time}
  H_u:=\inf\{n\geq0:S_n\geq u\}.
\end{equation}
The value of $H_u$ is allowed to be $+\infty$, although the strong law and
$\mu>0$ imply that it is finite almost surely.  For the ordinary real-valued
distributional statement below, define
\begin{equation}\label{eq:finite-hitting-version}
  \overline H_u
  :=
  \begin{cases}
    H_u,&H_u<\infty,\\
    0,&H_u=\infty.
  \end{cases}
\end{equation}
Thus $\overline H_u=H_u$ almost surely after identifying finite natural
numbers with real numbers.  The survival-probability statement does not
require this modification.

The form needed in Chapter 3 is the following survival-probability limit.

\begin{theorem}[First-passage central limit theorem]
\label{thm:first-passage-clt}
  For $a\in\R$ and $u$ sufficiently large, set
  \begin{equation}\label{eq:canonical-time}
    n_u(a)
    :=
    \left\lfloor
      \frac{u}{\mu}
      +a\frac{\sigma}{\mu^{3/2}}\sqrt{u}
    \right\rfloor.
  \end{equation}
  Then
  \begin{equation}\label{eq:first-passage-profile}
    \Pp\bigl(H_u>n_u(a)\bigr)
    \longrightarrow
    \Phi_{\mathrm G}(-a),
    \qquad u\longrightarrow\infty,
  \end{equation}
  where $\Phi_{\mathrm G}$ is the standard Gaussian distribution function.
  Equivalently,
  \begin{equation}\label{eq:first-passage-distribution}
    \frac{\overline H_u-u/\mu}
    {\sigma\mu^{-3/2}\sqrt{u}}
    \ind N(0,1).
  \end{equation}
\end{theorem}

We prove a slightly more flexible statement, which avoids manipulating the
extended-valued random variable $H_u$.

\begin{proposition}[Moving-level profile]\label{prop:moving-level-profile}
  Let $(u_r)_{r\geq1}$ be positive real numbers and let
  $(n_r)_{r\geq1}$ be positive natural numbers such that
  \begin{equation}\label{eq:moving-level-assumptions}
    n_r\longrightarrow\infty,
    \qquad
    c_r:=\frac{u_r-\mu n_r}{\sigma\sqrt{n_r}}
    \longrightarrow c\in\R.
  \end{equation}
  Then
  \begin{equation}\label{eq:moving-level-conclusion}
    \Pp(H_{u_r}>n_r)
    \longrightarrow
    \Phi_{\mathrm G}(c).
  \end{equation}
\end{proposition}

The proof has three parts.  First, reversal of a finite i.i.d. block
identifies the terminal drawdown $D_n$ with a running maximum of the
negative walk.  Second, the strong law makes these negative running maxima
uniformly tight.  Third, the ordinary central limit theorem and Slutsky's
theorem transfer the Gaussian limit from $S_n$ to $M_n$.

\section{Reversal and the terminal drawdown}\label{sec:drawdown}

For $0\leq j\leq n$, define the backward partial sums
\begin{equation}\label{eq:backward-sums}
  \overleftarrow S_j^{(n)}
  :=
  \sum_{r=0}^{j-1}\xi_{n-1-r},
  \qquad
  \overleftarrow S_0^{(n)}:=0,
\end{equation}
and define
\begin{equation}\label{eq:negative-maximum}
  Q_n:=\max_{0\leq j\leq n}(-S_j).
\end{equation}

\begin{lemma}[Finite reversal]\label{lem:finite-reversal}
  For every $n\geq0$,
  \begin{equation}\label{eq:drawdown-backward}
    D_n
    =
    \max_{0\leq j\leq n}\bigl(-\overleftarrow S_j^{(n)}\bigr)
  \end{equation}
  pointwise, and
  \begin{equation}\label{eq:drawdown-identically-distributed}
    D_n\stackrel{\mathrm d}=Q_n.
  \end{equation}
\end{lemma}

\begin{proof}
  By definition,
  \begin{align*}
    D_n
    &=\max_{0\leq k\leq n}(S_k-S_n)\\
    &=\max_{0\leq k\leq n}
      \left(-\sum_{i=k}^{n-1}\xi_i\right).
  \end{align*}
  Substituting $j=n-k$ gives \cref{eq:drawdown-backward}.

  The random vectors
  \begin{equation*}
    (\xi_{n-1},\xi_{n-2},\ldots,\xi_0)
    \quad\hbox{and}\quad
    (\xi_0,\xi_1,\ldots,\xi_{n-1})
  \end{equation*}
  have the same distribution.  Indeed, each vector has product law
  $\mathcal L(\xi_0)^{\otimes n}$; equivalently, the joint law of a finite
  i.i.d. family is invariant under every permutation of its indices.  The
  map which sends a vector $(x_0,\ldots,x_{n-1})$ to
  \begin{equation*}
    \max_{0\leq j\leq n}
    \left(-\sum_{r=0}^{j-1}x_r\right)
  \end{equation*}
  is continuous, hence measurable.  Applying this map to the two vectors
  proves \cref{eq:drawdown-identically-distributed}.
\end{proof}

The next lemma is the only place where positive drift is used in the
drawdown estimate.

\begin{lemma}[Uniform tightness of the drawdown]
\label{lem:drawdown-tight}
  For every $\varepsilon>0$, there exists $b<\infty$ such that
  \begin{equation}\label{eq:drawdown-tight}
    \sup_{n\geq0}\Pp(D_n>b)<\varepsilon.
  \end{equation}
\end{lemma}

\begin{proof}
  For $m\in\N$, let
  \begin{equation}\label{eq:lower-bound-event}
    A_m:=\{\omega:S_j(\omega)\geq-m\text{ for every }j\geq0\}.
  \end{equation}
  Each $A_m$ is measurable, and $A_m\subseteq A_{m+1}$.  The strong law gives
  \begin{equation}\label{eq:slln}
    \frac{S_j}{j}\longrightarrow\mu>0
    \qquad\hbox{almost surely}.
  \end{equation}
  On the event in \cref{eq:slln}, there is a random index $j_0$ such that
  $S_j\geq0$ for every $j\geq j_0$.  The finite set
  $\{S_0,\ldots,S_{j_0-1}\}$ has a finite minimum.  Hence there exists a
  natural number $m=m(\omega)$ such that $S_j(\omega)\geq-m$ for every
  $j\geq0$.  It follows that
  \begin{equation}\label{eq:lower-bound-exhaustion}
    \Pp\left(\bigcup_{m\geq0}A_m\right)=1.
  \end{equation}
  By continuity of probability from below,
  \begin{equation}\label{eq:lower-bound-probability}
    \Pp(A_m^c)\longrightarrow0.
  \end{equation}

  On $A_m$ we have $Q_n\leq m$ for every $n$.  By
  \cref{eq:drawdown-identically-distributed},
  \begin{equation}\label{eq:drawdown-tail-bound}
    \Pp(D_n>m)
    =\Pp(Q_n>m)
    \leq\Pp(A_m^c)
    \qquad\hbox{for every }n.
  \end{equation}
  Choose $m$ so that the right-hand side is smaller than $\varepsilon$.
  This proves \cref{eq:drawdown-tight}.
\end{proof}

\begin{corollary}[Negligible drawdown]
\label{cor:drawdown-negligible}
  If $b_n>0$ and $b_n\to\infty$, then
  \begin{equation}\label{eq:drawdown-in-probability}
    \frac{D_n}{b_n}\longrightarrow0
    \qquad\hbox{in probability}.
  \end{equation}
  In particular, as $n\to\infty$ through $n\geq1$,
  $D_n/(\sigma\sqrt n)\to0$ in probability.
\end{corollary}

\begin{proof}
  Fix $\delta,\varepsilon>0$.  By \cref{lem:drawdown-tight}, choose $m$ such
  that $\sup_n\Pp(D_n>m)<\varepsilon$.  For all sufficiently large $n$,
  $m<\delta b_n$, and therefore
  \begin{equation*}
    \Pp\left(\frac{D_n}{b_n}>\delta\right)
    \leq\Pp(D_n>m)
    <\varepsilon.
  \end{equation*}
\end{proof}

\begin{remark}\label{rem:why-reversal}
  No moment estimate for the all-time negative minimum is needed.  The strong
  law proves only that the walk is bounded below almost surely, and
  continuity from below converts this pathwise statement into the uniform
  tightness in \cref{lem:drawdown-tight}.  This avoids a maximal inequality,
  an overshoot estimate, and a functional central limit theorem.
\end{remark}

\section{Central limit theorem for the running maximum}
\label{sec:maximum-clt}

\begin{proposition}[CLT for the running maximum]
\label{prop:maximum-clt}
  Under \cref{eq:moments},
  \begin{equation}\label{eq:maximum-clt}
    \frac{M_n-\mu n}{\sigma\sqrt n}
    \ind N(0,1).
  \end{equation}
\end{proposition}

\begin{proof}
  The ordinary central limit theorem gives
  \begin{equation}\label{eq:sum-clt}
    Z_n:=\frac{S_n-\mu n}{\sigma\sqrt n}
    \ind Z,
    \qquad Z\sim N(0,1).
  \end{equation}
  By \cref{eq:walk-maximum-drawdown},
  \begin{equation}\label{eq:maximum-decomposition}
    \frac{M_n-\mu n}{\sigma\sqrt n}
    =
    Z_n+\frac{D_n}{\sigma\sqrt n}.
  \end{equation}
  The second term converges to zero in probability by
  \cref{cor:drawdown-negligible}.  Slutsky's theorem applied to
  \cref{eq:maximum-decomposition} proves \cref{eq:maximum-clt}.
\end{proof}

We now pass from the maximum to the passage event.  The precise endpoint is
important:
\begin{equation}\label{eq:hitting-maximum-event}
  \{H_u>n\}
  =
  \{S_k<u\text{ for every }0\leq k\leq n\}
  =
  \{M_n<u\}.
\end{equation}
The strict inequality agrees with the weak inequality in the definition of
$H_u$.

\begin{proof}[Proof of \cref{prop:moving-level-profile}]
  Define
  \begin{equation*}
    W_n:=\frac{M_n-\mu n}{\sigma\sqrt n}.
  \end{equation*}
  Since $n_r\to\infty$, \cref{prop:maximum-clt} along the subsequence
  $(n_r)$ gives $W_{n_r}\ind Z$.  The deterministic variables $-c_r$
  converge in probability to $-c$.  Slutsky's theorem therefore gives
  \begin{equation}\label{eq:shifted-maximum-clt}
    W_{n_r}-c_r\ind Z-c.
  \end{equation}
  The boundary of $(-\infty,0)$ is the singleton $\{0\}$, and the law of
  $Z-c$ has no atom there.  Portmanteau applied to
  \cref{eq:shifted-maximum-clt} yields
  \begin{equation*}
    \Pp(W_{n_r}-c_r<0)
    \longrightarrow
    \Pp(Z-c<0)
    =\Phi_{\mathrm G}(c).
  \end{equation*}
  By the definition of $c_r$ and \cref{eq:hitting-maximum-event},
  \begin{equation*}
    \{W_{n_r}-c_r<0\}
    =\{M_{n_r}<u_r\}
    =\{H_{u_r}>n_r\}.
  \end{equation*}
  This proves \cref{eq:moving-level-conclusion}.
\end{proof}

\section{Floors, moving levels, and moving times}
\label{sec:moving-parameters}

The remaining work is deterministic.  We state it in the sequential form
used by a Lean proof; the same calculation holds along the filter $u\to
\infty$.

\begin{lemma}[Canonical floor asymptotics]
\label{lem:canonical-floor}
  Let $L_r\to\infty$, let $a\in\R$, and let $q_r=o(\sqrt{L_r})$.  For all
  sufficiently large $r$, define
  \begin{equation}\label{eq:perturbed-time}
    n_r
    :=
    \left\lfloor
      \frac{L_r}{\mu}
      +a\frac{\sigma}{\mu^{3/2}}\sqrt{L_r}
      +q_r
    \right\rfloor.
  \end{equation}
  If
  \begin{equation}\label{eq:perturbed-level}
    \widetilde L_r=L_r+o(\sqrt{L_r}),
  \end{equation}
  then
  \begin{equation}\label{eq:floor-asymptotics}
    n_r\longrightarrow\infty,
    \qquad
    \frac{n_r}{L_r}\longrightarrow\frac1\mu,
    \qquad
    \frac{\widetilde L_r-\mu n_r}{\sigma\sqrt{n_r}}
    \longrightarrow-a.
  \end{equation}
\end{lemma}

\begin{proof}
  Write
  \begin{equation}\label{eq:floor-error}
    n_r
    =
    \frac{L_r}{\mu}
    +a\frac{\sigma}{\mu^{3/2}}\sqrt{L_r}
    +q_r+e_r,
    \qquad -1<e_r\leq0.
  \end{equation}
  Dividing by $L_r$ and using $q_r=o(\sqrt{L_r})$ proves
  $n_r/L_r\to1/\mu$, and in particular $n_r\to\infty$.  It also gives
  \begin{equation}\label{eq:sqrt-time-asymptotic}
    \frac{\sqrt{n_r}}{\sqrt{L_r}}
    \longrightarrow\frac1{\sqrt\mu}.
  \end{equation}
  From \cref{eq:floor-error},
  \begin{equation}\label{eq:center-floor-expansion}
    \widetilde L_r-\mu n_r
    =
    (\widetilde L_r-L_r)
    -a\frac{\sigma}{\sqrt\mu}\sqrt{L_r}
    -\mu q_r-\mu e_r.
  \end{equation}
  Divide \cref{eq:center-floor-expansion} by
  $\sigma\sqrt{n_r}$.  The first, third, and fourth terms tend to zero by
  \cref{eq:perturbed-level}, the hypothesis on $q_r$, and the bound on
  $e_r$.  The second tends to $-a$ by
  \cref{eq:sqrt-time-asymptotic}.
\end{proof}

\begin{corollary}[Moving-level and time profile]
\label{cor:moving-level-time-profile}
  Under the hypotheses of \cref{lem:canonical-floor},
  \begin{equation}\label{eq:moving-level-time-profile}
    \Pp(H_{\widetilde L_r}>n_r)
    \longrightarrow
    \Phi_{\mathrm G}(-a).
  \end{equation}
\end{corollary}

\begin{proof}
  Apply \cref{prop:moving-level-profile} with $u_r=\widetilde L_r$ and
  $c=-a$, using \cref{eq:floor-asymptotics}.
\end{proof}

The paper also perturbs a time which has already been rounded down.  We record
that convention separately.

\begin{corollary}[Perturbation after flooring]
\label{cor:post-floor-profile}
  Fix $a\in\R$.  Let $L_r\to\infty$, let
  $\widetilde L_r=L_r+o(\sqrt{L_r})$, and let
  $v_r=o(\sqrt{L_r})$.  Define
  \begin{align}
    t_r(a)
    &:=
    \left\lfloor
      \frac{L_r}{\mu}
      +a\frac{\sigma}{\mu^{3/2}}\sqrt{L_r}
    \right\rfloor,\label{eq:base-floored-time}\\
    \widehat n_r
    &:=\left\lfloor t_r(a)+v_r\right\rfloor.
    \label{eq:post-floor-time}
  \end{align}
  For all sufficiently large $r$, both quantities are positive and are
  regarded as natural-valued times.  Then
  \begin{equation}\label{eq:post-floor-profile}
    \Pp(H_{\widetilde L_r}>\widehat n_r)
    \longrightarrow
    \Phi_{\mathrm G}(-a).
  \end{equation}
\end{corollary}

\begin{proof}
  Put
  \begin{equation*}
    x_r
    :=
    \frac{L_r}{\mu}
    +a\frac{\sigma}{\mu^{3/2}}\sqrt{L_r},
    \qquad
    n_r':=\lfloor x_r+v_r\rfloor.
  \end{equation*}
  Since $t_r(a)=\lfloor x_r\rfloor$ is an integer,
  \begin{equation}\label{eq:two-floor-error}
    \abs{\widehat n_r-n_r'}\leq1.
  \end{equation}
  Indeed, writing $x_r=\lfloor x_r\rfloor+\{x_r\}$ with
  $0\leq\{x_r\}<1$, the two floors differ by
  $\lfloor v_r+\{x_r\}\rfloor-\lfloor v_r\rfloor\in\{0,1\}$.
  By \cref{lem:canonical-floor} with $q_r=v_r$,
  \begin{equation*}
    n_r'\longrightarrow\infty,
    \qquad
    \frac{\widetilde L_r-\mu n_r'}{\sigma\sqrt{n_r'}}
    \longrightarrow-a.
  \end{equation*}
  The bounded difference in \cref{eq:two-floor-error} implies the same two
  limits with $\widehat n_r$ in place of $n_r'$.  The conclusion now follows
  from \cref{prop:moving-level-profile}.
\end{proof}

\begin{proof}[Proof of \cref{thm:first-passage-clt}]
  Let $(u_r)$ be an arbitrary sequence with $u_r\to\infty$.  Take
  $L_r=\widetilde L_r=u_r$ and $q_r=0$ in
  \cref{cor:moving-level-time-profile}.  This proves
  \cref{eq:first-passage-profile} along every sequence tending to infinity.
  The sequential characterization of limits on $\R$ gives the asserted
  limit as $u\to\infty$.

  Since $H_u$ is integer-valued and finite almost surely, the two events
  below agree up to a null set:
  \begin{equation*}
    \1_{\left\{
        (\overline H_u-u/\mu)/(\sigma\mu^{-3/2}\sqrt u)>a
      \right\}}
    =
    \1_{\{H_u>n_u(a)\}}
    \qquad\hbox{almost surely}.
  \end{equation*}
  Thus \cref{eq:first-passage-profile} gives convergence of the survival
  functions at every $a$.  The identity
  $1-\Phi_{\mathrm G}(-a)=\Phi_{\mathrm G}(a)$ and continuity of the Gaussian
  distribution function give \cref{eq:first-passage-distribution}.
\end{proof}

\section{An almost surely bounded increasing correction}
\label{sec:bounded-correction}

The Chapter 3 logarithmic radius is not exactly $S_n$.  It is $S_n+E_n$,
where $E_n$ is increasing and converges almost surely to a finite random
variable.  No independence between $E$ and the increments is needed.

Let $(E_n)_{n\geq0}$ and $E_\infty$ be random variables such that
\begin{equation}\label{eq:correction-assumptions}
  \begin{split}
    \Omega_E:=\{\omega\in\Omega:\;&
      0\leq E_n(\omega)\leq E_{n+1}(\omega)
        \leq E_\infty(\omega)<\infty\\
      &\text{for every }n\geq0\},
      \qquad \Pp(\Omega_E)=1.
  \end{split}
\end{equation}
For $u>0$, set
\begin{equation}\label{eq:corrected-hitting-time}
  T_u:=\inf\{n\geq0:S_n+E_n\geq u\}.
\end{equation}

\begin{proposition}[Passage with a bounded increasing correction]
\label{prop:corrected-passage}
  Fix $a\in\R$.  Suppose $L_r\to\infty$,
  $\widetilde L_r=L_r+o(\sqrt{L_r})$, and
  $q_r=o(\sqrt{L_r})$, and let $n_r$ be given by
  \cref{eq:perturbed-time}.  Then
  \begin{equation}\label{eq:corrected-profile}
    \Pp(T_{\widetilde L_r}>n_r)
    \longrightarrow
    \Phi_{\mathrm G}(-a).
  \end{equation}
\end{proposition}

\begin{proof}
  Fix $B\geq0$.  For every $u>B$, on the event
  $\Omega_E\cap\{E_\infty\leq B\}$ we have
  \begin{equation}\label{eq:corrected-sandwich}
    H_{u-B}\leq T_u\leq H_u.
  \end{equation}
  Indeed, $S_n\geq u$ implies $S_n+E_n\geq u$, while
  $S_n+E_n\geq u$ and $E_n\leq B$ imply $S_n\geq u-B$.
  Since $\Pp(\Omega_E^c)=0$, the exceptional event in the resulting
  probability bounds may be taken to be $\{E_\infty>B\}$ alone.
  Since $\widetilde L_r\to\infty$, we may take
  $u=\widetilde L_r>B$ for all sufficiently large $r$.  Consequently,
  \begin{align}\label{eq:corrected-probability-sandwich}
    \Pp(H_{\widetilde L_r-B}>n_r)-\Pp(E_\infty>B)
    &\leq \Pp(T_{\widetilde L_r}>n_r)\notag\\
    &\leq
    \Pp(H_{\widetilde L_r}>n_r)+\Pp(E_\infty>B).
  \end{align}
  Both uncorrected probabilities converge to
  $\Phi_{\mathrm G}(-a)$ by
  \cref{cor:moving-level-time-profile}, because subtracting the fixed number
  $B$ is an $o(\sqrt{L_r})$ perturbation of the level.  Taking lower and upper
  limits in \cref{eq:corrected-probability-sandwich} gives an error at most
  $\Pp(E_\infty>B)$.  Finally,
  \begin{equation*}
    \Pp(E_\infty>B)\longrightarrow0,
    \qquad B\longrightarrow\infty,
  \end{equation*}
  since $E_\infty<\infty$ almost surely.  This proves
  \cref{eq:corrected-profile}.
\end{proof}

\begin{corollary}[Corrected perturbation after flooring]
\label{cor:corrected-post-floor-profile}
  Fix $a\in\R$.  Under \cref{eq:correction-assumptions}, let $L_r$,
  $\widetilde L_r$, $v_r$, and $\widehat n_r$ satisfy the hypotheses and definitions of
  \cref{cor:post-floor-profile}.  Then
  \begin{equation}\label{eq:corrected-post-floor-profile}
    \Pp(T_{\widetilde L_r}>\widehat n_r)
    \longrightarrow
    \Phi_{\mathrm G}(-a).
  \end{equation}
\end{corollary}

\begin{proof}
  Fix $B\geq0$.  For all sufficiently large $r$, the sandwich
  \cref{eq:corrected-sandwich} and the same event argument as in
  \cref{eq:corrected-probability-sandwich} give
  \begin{align*}
    \Pp(H_{\widetilde L_r-B}>\widehat n_r)-\Pp(E_\infty>B)
    &\leq \Pp(T_{\widetilde L_r}>\widehat n_r)\\
    &\leq
    \Pp(H_{\widetilde L_r}>\widehat n_r)+\Pp(E_\infty>B).
  \end{align*}
  Both uncorrected probabilities converge to
  $\Phi_{\mathrm G}(-a)$ by \cref{cor:post-floor-profile}; the fixed shift
  $-B$ preserves the hypothesis
  $\widetilde L_r=L_r+o(\sqrt{L_r})$.  Taking lower and upper limits and then
  sending $B\to\infty$ proves the result.
\end{proof}

\section{Application to the fixed-width Gaussian radius}
\label{sec:gaussian-application}

Fix $N\geq1$ and $0<A<A_c(N)$.  Let $(\chi_{N,j})_{j\geq0}$ be independent
copies of a chi random variable with $N$ degrees of freedom, and put
\begin{equation}\label{eq:gaussian-increments}
  \begin{aligned}
    \xi_j
    &:=-\log\left(\frac{A}{\sqrt N}\chi_{N,j}\right),\\
    \mu&=\gamma_{A,N}>0,\qquad
    \sigma_N:=\sqrt{\Var(\log\chi_N)}>0,\qquad
    \sigma=\sigma_N.
  \end{aligned}
\end{equation}
This is the zero-based reindexing of the paper's sequence
$(\chi_{N,j})_{j\geq1}$ and matches \cref{eq:walk-maximum-drawdown}.
On a single event of probability one, the logarithmic radius reduction in
Chapter 3 gives, simultaneously for every $t\in\N$,
\begin{equation}\label{eq:log-radius-reduction}
  -\log R_t
  =
  -\log R_0+S_t+E_t,
\end{equation}
where the correction satisfies \cref{eq:correction-assumptions}.
The application uses the following three inputs from that reduction.  The
Gaussian innovations, and hence the variables $\chi_{N,j}$, are independent;
$\log\chi_N$ is square-integrable and nonconstant; and $R_t>0$ for every
$t\in\N$ on one event of probability one.  The last assertion follows by
taking the countable intersection of the fixed-time nonvanishing events.  On
this common event, taking logarithms in the definition of the radius entrance
time is legitimate.
For a radius threshold $\varepsilon_\varrho$, set
\begin{equation}\label{eq:paper-levels}
  L_\varrho:=\log\frac{R_0}{\varrho},
  \qquad
  L_{\varepsilon_\varrho}:=
  \log\frac{R_0}{\varepsilon_\varrho}.
\end{equation}
Then the entrance time
\begin{equation*}
  \theta_{\varepsilon_\varrho}^{\mathrm{rad}}
  :=\inf\{t\geq0:R_t\leq\varepsilon_\varrho\}
\end{equation*}
agrees almost surely with $T_{L_{\varepsilon_\varrho}}$ from
\cref{eq:corrected-hitting-time}.  In particular, their survival
probabilities at every deterministic time are equal.

Assume
\begin{equation}\label{eq:paper-moving-level}
  L_{\varepsilon_\varrho}
  =L_\varrho+o(\sqrt{L_\varrho}),
  \qquad \varrho\downarrow0.
\end{equation}
To justify the limits below from the sequence-indexed results, fix an
arbitrary sequence of positive meshes $\varrho_r\to0$.  Apply
\cref{prop:corrected-passage,cor:corrected-post-floor-profile} with
$L_r=L_{\varrho_r}$ and
$\widetilde L_r=L_{\varepsilon_{\varrho_r}}$, and with
$v_r=u_{\varrho_r}$ for the post-floor assertion.  Since the mesh sequence is
arbitrary, the sequential characterization of limits on $\R$ gives the
stated limits as $\varrho\downarrow0$.

For every fixed $a\in\R$, since $L_\varrho\to\infty$,
\cref{prop:corrected-passage} with
$\mu=\gamma_{A,N}$ and $\sigma=\sigma_N$ gives, for
\begin{equation}\label{eq:paper-cutoff-time}
  t_\varrho(a)
  :=
  \left\lfloor
    \frac{L_\varrho}{\gamma_{A,N}}
    +a\frac{\sigma_N}{\gamma_{A,N}^{3/2}}
      \sqrt{L_\varrho}
  \right\rfloor,
\end{equation}
the limit
\begin{equation}\label{eq:paper-profile}
  \Pp\left(
    \theta_{\varepsilon_\varrho}^{\mathrm{rad}}>t_\varrho(a)
  \right)
  \longrightarrow
  \Phi_{\mathrm G}(-a).
\end{equation}
If the time in \cref{eq:paper-cutoff-time} is perturbed by a deterministic
$u_\varrho=o(\sqrt{L_\varrho})$ term after the first floor, then
\cref{cor:corrected-post-floor-profile} gives
\begin{equation*}
  \Pp\left(
    \theta_{\varepsilon_\varrho}^{\mathrm{rad}}
    >\left\lfloor t_\varrho(a)+u_\varrho\right\rfloor
  \right)
  \longrightarrow
  \Phi_{\mathrm G}(-a).
\end{equation*}
This is the exact convention in the paper's Gaussian moving-threshold
proposition.  No separate renewal theorem or external stopped-random-walk
theorem is needed.

\section{The formalized statements}\label{sec:formalized}

The proof above is carried out in \texttt{AbsorptionCutoff/FirstPassageCLT.lean}, a
pure-Mathlib leaf module: it imports nothing else from the development, so the
argument stands independently of the rounded chain.  The correspondence is the
following.

The declaration-by-declaration map is given in
\texttt{notes/formalized-libraries.tex}, which covers this result together with
the key renewal theorem.

The external inputs are the scalar i.i.d.\ central limit theorem
\lean{tendstoInDistribution\_inv\_sqrt\_mul\_sum\_sub}, the strong law
\lean{strong\_law\_ae}, the identification of the laws of the forward and
reversed blocks through \texttt{IdentDistrib.pi}, the two Slutsky steps, and
the Portmanteau theorem in the form
\lean{tendsto\_measure\_of\_null\_frontier\_of\_tendsto'}, applied to
$(-\infty,0)$, whose frontier $\{0\}$ is null for a nondegenerate Gaussian
law.  The hitting time is Mathlib's \texttt{hittingAfter}.

Two endpoint conventions are used throughout.  The passage time $H_u$ hits
the closed set $[u,\infty)$, so survival through time $n$ is the strict event
$M_n<u$; and the floor error in \cref{eq:floor-error} lies in $(-1,0]$.
Neither affects the Gaussian limit.

\end{document}
